% TeX root=../main.tex

\begin{table}[!htb]
	\begin{adjustwidth}{0cm}{0cm}
		\caption[Temas neutros em oclusiva e em \emph{-s-}]{Temas neutros em oclusiva e em \emph{-s-} (Paradigmas:
			\ocsex{agnẽ} `cordeiro', \ocsex{slovo} `palavra')}\label{tab:temasemoclusivaes}
		\begin{center}
			\begin{tabular}[c]{llllll}
				\toprule
				                   &                  & \multicolumn{2}{c}{neutros em \emph{-t-}} & \multicolumn{2}{c}{neutros em \emph{-s-}}                                            \\
				\midrule
				\emph{sg.}         & \emph{nom.acu}   & \cirilex{agnẽ}                            & \ocsex{agnẽ}                              & \cirilex{slovo}     & \ocsex{slovo}      \\
				                   &                  &                                           &                                           & \cirilex{slovesI}   & \ocsex{sloves-I}   \\
				\rowcolor{gray!10} & \emph{gen.}      & \cirilex{agnẽte}                          & \ocsex{agnẽt-e}                           & \cirilex{slovese}   & \ocsex{sloves-e}   \\
				\rowcolor{gray!10} &                  &                                           &                                           & \cirilex{slovesi}   & \ocsex{sloves-i}   \\
				                   & \emph{dat.}      & \cirilex{agnẽti}                          & \ocsex{agnẽt-i}                           & \cirilex{slovesu}   & \ocsex{sloves-u}   \\
				                   & \emph{inst.}     & \cirilex{agnẽtemI}                        & \ocsex{agnẽt-emI}                         & \cirilex{slovesemI} & \ocsex{sloves-emI} \\
				                   &                  &                                           &                                           & \cirilex{slovesImI} & \ocsex{sloves-ImI} \\
				\rowcolor{gray!10} & \emph{loc.}      & \cirilex{agnẽte}                          & \ocsex{agnẽt-e}                           & \cirilex{slovese}   & \ocsex{sloves-e}   \\
				\rowcolor{gray!10} &                  & \cirilex{agnẽti}                          & \ocsex{agnẽt-i}                           & \cirilex{slovesi}   & \ocsex{sloves-i}   \\
				\midrule
				\emph{du.}         & \emph{nom.acu}   & \cirilex{agnẽtE}                          & \ocsex{agnẽt-E}                           & \cirilex{slovesE}   & \ocsex{sloves-E}   \\
				                   &                  &                                           &                                           & \cirilex{slovesi}   & \ocsex{sloves-i}   \\
				                   &                  &                                           &                                           &                     &                    \\
				\rowcolor{gray!10} & \emph{gen.loc.}  & \cirilex{agnẽtu}                          & \ocsex{agnẽt-u}                           & \cirilex{slovesu}   & \ocsex{sloves-u}   \\
				                   & \emph{dat.inst.} & \cirilex{agnẽtIma}                        & \ocsex{agnẽt-Ima}                         & \cirilex{slovosIma} & \ocsex{sloves-Ima} \\
				\midrule
				\emph{pl.}         & \emph{nom.acu}   & \cirilex{agnẽta}                          & \ocsex{agnẽt-a}                           & \cirilex{slovesa}   & \ocsex{sloves-a}   \\
				\rowcolor{gray!10} & \emph{gen.}      & \cirilex{agnẽtU}                          & \ocsex{agnẽt-U}                           & \cirilex{slovesU}   & \ocsex{sloves-U}   \\
				                   & \emph{dat.}      & \cirilex{agnẽtemU}                        & \ocsex{agnẽt-emU}                         & \cirilex{slovesemU} & \ocsex{sloves-emU} \\
				\rowcolor{gray!10} & \emph{inst.}     & \cirilex{agnẽty}                          & \ocsex{agnẽt-y}                           & \cirilex{slovesy}   & \ocsex{sloves-y}   \\
				                   & \emph{loc.}      & \cirilex{agnẽtexU}                        & \ocsex{agnẽt-exU}                         & \cirilex{slovesExU} & \ocsex{sloves-ExU} \\
				\bottomrule
			\end{tabular}
		\end{center}
	\end{adjustwidth}
\end{table}
